\documentclass[preview]{standalone}

\usepackage{amsmath}
\usepackage{amssymb}
\usepackage{stellar}
\usepackage{definitions}
\usepackage{bettelini}

\begin{document}

\id{linearalgebra-linear-isometries}
\genpage

\section{Linear isometries}

\begin{snippetdefinition}{linear-isometry-definition}{Linear isometry}
    Let \((V,\langle\cdot,\cdot\rangle)\) be an \innerproductspace over a
    \field \(\mathbb{F}\in\{\realnumbers,\complexnumbers\}\). An
    endomorphism \(f\in\setoflineartransformations(V,V)\) is a
    \emph{linear isometry} if
    \[
        \langle f(v),f(w)\rangle=\langle v,w\rangle
    \]
    for every \(v,w\in V\).
\end{snippetdefinition}

\begin{snippetproposition}{linear-isometry-preserves-norms}{Linear isometries preserve norms}
    Let \((V,\langle\cdot,\cdot\rangle)\) be an \innerproductspace over a
    \field \(\mathbb{F}\in\{\realnumbers,\complexnumbers\}\). If
    \(f\in\setoflineartransformations(V,V)\) is a \linearisometry, then
    \[
        \|f(v)\|=\|v\|
    \]
    for every \(v\in V\).
\end{snippetproposition}

\begin{snippetproof}{linear-isometry-preserves-norms-proof}{linear-isometry-preserves-norms}{Linear isometries preserve norms}
    For every \(v\in V\),
    \[
        \|f(v)\|^2
        =
        \langle f(v),f(v)\rangle
        =
        \langle v,v\rangle
        =
        \|v\|^2.
    \]
    Since \norm[norms] are non-negative, \(\|f(v)\|=\|v\|\).
\end{snippetproof}

\begin{snippetexample}{complex-line-linear-isometries}{Linear isometries of the complex line}
    Consider \(\complexnumbers\) as a complex \innerproductspace with the
    canonical \innerproduct. If
    \(f\in\setoflineartransformations(\complexnumbers,\complexnumbers)\),
    then there exists \(\lambda\in\complexnumbers\) such that
    \[
        f(z)=\lambda z
    \]
    for every \(z\in\complexnumbers\). The map \(f\) is a
    \linearisometry if and only if
    \[
        |\lambda|=1.
    \]
    Indeed,
    \[
        \|f(z)\|=|\lambda z|=|\lambda|\,|z|.
    \]
    Hence \(f\) preserves \norm[norms] for every \(z\in\complexnumbers\) if and
    only if \(|\lambda|=1\). In this case
    \[
        \lambda=\cos\theta+\mathrm{i}\sin\theta=e^{\mathrm{i}\theta}
    \]
    for some \(\theta\in\realnumbers\).
\end{snippetexample}

\begin{snippetexample}{plane-rotation-linear-isometry-example}{Plane rotations are linear isometries}
    For every \(\theta\in\realnumbers\), let
    \[
        R_\theta=
        \begin{pmatrix}
            \cos\theta & -\sin\theta\\
            \sin\theta & \cos\theta
        \end{pmatrix}.
    \]
    The associated map
    \[
        L_{R_\theta}\colon\realnumbers^2\fromto\realnumbers^2,
        \qquad
        L_{R_\theta}(x)=R_\theta x
    \]
    is a \linearisometry for the canonical \innerproduct on
    \(\realnumbers^2\). Indeed, the columns of \(R_\theta\) are
    \orthonormal and
    \[
        R_\theta^\transpose R_\theta=\identmatrix{2}.
    \]
    Geometrically, \(L_{R_\theta}\) is the counterclockwise rotation of
    angle \(\theta\).
\end{snippetexample}

\section{Characterizations}

\begin{snippettheorem}{linear-isometry-characterization}{Characterization of linear isometries}
    Let \((V,\langle\cdot,\cdot\rangle)\) be a finite-dimensional
    \innerproductspace over a
    \field \(\mathbb{F}\in\{\realnumbers,\complexnumbers\}\), and let
    \(f\in\setoflineartransformations(V,V)\). The following are
    equivalent:
    \begin{enumerate}
        \item \(f\) is a \linearisometry;
        \item \(\|f(v)\|=\|v\|\) for every \(v\in V\);
        \item for every \orthonormalbasis
            \(\mathcal{B}=\{e_1,\ldots,e_n\}\) of \(V\), the collection
            \(\{f(e_1),\ldots,f(e_n)\}\) is an \orthonormalbasis of \(V\);
        \item there exists an \orthonormalbasis
            \(\mathcal{B}=\{e_1,\ldots,e_n\}\) of \(V\) such that
            \(\{f(e_1),\ldots,f(e_n)\}\) is an \orthonormalbasis of \(V\);
        \item \(f^\ast\circ f=\identityfunc_V\);
        \item \(f\circ f^\ast=\identityfunc_V\);
        \item \(f^\ast\) is a \linearisometry;
        \item \(f\) is invertible and \(f^{-1}=f^\ast\).
    \end{enumerate}
\end{snippettheorem}

\begin{snippetproof}{linear-isometry-characterization-proof}{linear-isometry-characterization}{Characterization of linear isometries}
    The implication \((1)\Rightarrow(2)\) follows from
    \[
        \|f(v)\|^2
        =
        \langle f(v),f(v)\rangle
        =
        \langle v,v\rangle
        =
        \|v\|^2.
    \]

    We prove \((2)\Rightarrow(1)\). If \(\mathbb{F}=\realnumbers\), the
    polarization identity gives
    \[
        \langle x,y\rangle
        =
        \frac{\|x+y\|^2-\|x-y\|^2}{4}.
    \]
    Since \(f\) is linear and preserves \norm[norms],
    \[
        \langle f(v),f(w)\rangle
        =
        \frac{\|f(v+w)\|^2-\|f(v-w)\|^2}{4}
        =
        \langle v,w\rangle.
    \]
    If \(\mathbb{F}=\complexnumbers\), the complex polarization identity is
    \[
        \langle x,y\rangle
        =
        \frac{\|x+y\|^2-\|x-y\|^2}{4}
        +
        \frac{\mathrm{i}\bigl(\|x+\mathrm{i}y\|^2-\|x-\mathrm{i}y\|^2\bigr)}{4}.
    \]
    Applying it to \(f(v)\) and \(f(w)\) gives again
    \(\langle f(v),f(w)\rangle=\langle v,w\rangle\).

    We prove \((1)\Rightarrow(3)\). Let
    \[
        \mathcal{B}=\{e_1,\ldots,e_n\}
    \]
    be an \orthonormalbasis of \(V\). Then
    \[
        \langle f(e_i),f(e_j)\rangle
        =
        \langle e_i,e_j\rangle
        =
        \delta_{ij}.
    \]
    Thus \(\{f(e_1),\ldots,f(e_n)\}\) is an \orthonormal collection of
    \(n=\lineardim V\) vectors, and therefore an \orthonormalbasis of
    \(V\).

    The implication \((3)\Rightarrow(4)\) is immediate.

    We prove \((4)\Rightarrow(5)\). Let
    \[
        \mathcal{B}=\{e_1,\ldots,e_n\}
    \]
    be an \orthonormalbasis such that
    \(\{f(e_1),\ldots,f(e_n)\}\) is an \orthonormalbasis. For every
    \(i,j\),
    \[
        \langle (f^\ast\circ f)(e_i),e_j\rangle
        =
        \langle f(e_i),f(e_j)\rangle
        =
        \delta_{ij}
        =
        \langle e_i,e_j\rangle.
    \]
    Since this holds for every \(j\), \((f^\ast\circ f)(e_i)=e_i\) for
    every \(i\), and hence \(f^\ast\circ f=\identityfunc_V\).

    We prove \((5)\Rightarrow(6)\). If
    \(f^\ast\circ f=\identityfunc_V\), then \(f\) is injective. Since
    \(V\) is finite-dimensional and \(f\) is an endomorphism, \(f\) is
    invertible. The equality \(f^\ast\circ f=\identityfunc_V\) gives
    \(f^\ast=f^{-1}\), and therefore
    \[
        f\circ f^\ast=\identityfunc_V.
    \]

    We prove \((6)\Rightarrow(7)\). For every \(v,w\in V\),
    \[
        \langle f^\ast(v),f^\ast(w)\rangle
        =
        \langle v,(f\circ f^\ast)(w)\rangle
        =
        \langle v,w\rangle.
    \]
    Hence \(f^\ast\) is a \linearisometry.

    We prove \((7)\Rightarrow(8)\). Since \(f^\ast\) is a
    \linearisometry, the implication \((1)\Rightarrow(5)\) applied to
    \(f^\ast\) gives
    \[
        (f^\ast)^\ast\circ f^\ast=\identityfunc_V.
    \]
    Therefore \(f\circ f^\ast=\identityfunc_V\). In finite dimension this
    implies that \(f\) is invertible and \(f^{-1}=f^\ast\).

    Finally, \((8)\Rightarrow(1)\). If \(f^{-1}=f^\ast\), then
    \(f^\ast\circ f=\identityfunc_V\). For every \(v,w\in V\),
    \[
        \langle f(v),f(w)\rangle
        =
        \langle v,(f^\ast\circ f)(w)\rangle
        =
        \langle v,w\rangle.
    \]
    Thus \(f\) is a \linearisometry.
\end{snippetproof}

\begin{snippetcorollary}{linear-isometries-are-normal}{Linear isometries are normal}
    Let \((V,\langle\cdot,\cdot\rangle)\) be a finite-dimensional
    \innerproductspace over a
    \field \(\mathbb{F}\in\{\realnumbers,\complexnumbers\}\). If
    \(f\in\setoflineartransformations(V,V)\) is a \linearisometry, then
    \(f\) is \normalendomorphism[normal].
\end{snippetcorollary}

\begin{snippetproof}{linear-isometries-are-normal-proof}{linear-isometries-are-normal}{Linear isometries are normal}
    By the characterization of \linearisometry[linear isometries],
    \(f^{-1}=f^\ast\). Hence
    \[
        f^\ast\circ f=f\circ f^\ast=\identityfunc_V.
    \]
    In particular \(f\circ f^\ast=f^\ast\circ f\), so \(f\) is
    \normalendomorphism[normal].
\end{snippetproof}

\section{Spectral characterization}

\begin{snippettheorem}{complex-linear-isometry-spectral-characterization}{Spectral characterization of complex linear isometries}
    Let \((V,\langle\cdot,\cdot\rangle)\) be a finite-dimensional complex
    \innerproductspace, and let
    \(f\in\setoflineartransformations(V,V)\). The following are
    equivalent:
    \begin{enumerate}
        \item \(f\) is a \linearisometry;
        \item there exists an \orthonormalbasis of \(V\) made of
            eigenvectors of \(f\), and every corresponding eigenvalue has
            modulus \(1\).
    \end{enumerate}
\end{snippettheorem}

\begin{snippetproof}{complex-linear-isometry-spectral-characterization-proof}{complex-linear-isometry-spectral-characterization}{Spectral characterization of complex linear isometries}
    \begin{iffproof}
        \item Suppose that \(f\) is a \linearisometry. Then
            \(f^{-1}=f^\ast\), so
            \[
                f^\ast\circ f=f\circ f^\ast=\identityfunc_V.
            \]
            Hence \(f\) is \normalendomorphism[normal]. By the complex
            \spectraltheorem, \(V\) has an \orthonormalbasis
            \[
                \mathcal{B}=\{e_1,\ldots,e_n\}
            \]
            made of eigenvectors of \(f\). Write
            \[
                f(e_i)=\lambda_i e_i.
            \]
            Since \(f\) is a \linearisometry,
            \[
                1=\|e_i\|=\|f(e_i)\|=|\lambda_i|\,\|e_i\|=|\lambda_i|.
            \]

        \item Conversely, suppose that \(V\) has an \orthonormalbasis
            \[
                \mathcal{B}=\{e_1,\ldots,e_n\}
            \]
            such that
            \[
                f(e_i)=\lambda_i e_i,
                \qquad
                |\lambda_i|=1
            \]
            for every \(i\). Let \(v\in V\). By the \parsevalidentity,
            \[
                v=\langle v,e_1\rangle e_1+\cdots+\langle v,e_n\rangle e_n.
            \]
            Therefore
            \[
                f(v)=
                \lambda_1\langle v,e_1\rangle e_1+\cdots+
                \lambda_n\langle v,e_n\rangle e_n.
            \]
            Again by the \parsevalidentity,
            \[
                \|f(v)\|^2
                =
                |\lambda_1|^2|\langle v,e_1\rangle|^2+\cdots+
                |\lambda_n|^2|\langle v,e_n\rangle|^2
                =
                \|v\|^2.
            \]
            Hence \(f\) preserves \norm[norms]. By the characterization of
            \linearisometry[linear isometries], \(f\) is a
            \linearisometry.
    \end{iffproof}
\end{snippetproof}

\section{Matrix form}

\begin{snippet}{linear-isometries-matrix-setup}
    Let \(\mathbb{F}\in\{\realnumbers,\complexnumbers\}\). Consider
    \(\mathbb{F}^n\) with the canonical \innerproduct. If
    \(A\in\matrices_{n\times n}(\mathbb{F})\), denote by
    \[
        L_A\colon\mathbb{F}^n\fromto\mathbb{F}^n,
        \qquad
        L_A(x)=Ax
    \]
    the associated \lineartransformation.
\end{snippet}

\begin{snippetproposition}{matrix-linear-isometry-criterion}{Matrix criterion for linear isometries}
    Let \(\mathbb{F}\in\{\realnumbers,\complexnumbers\}\), let
    \(\mathbb{F}^n\) be endowed with the canonical \innerproduct, and let
    \(A\in\matrices_{n\times n}(\mathbb{F})\). The associated map
    \(L_A\) is a \linearisometry if and only if \(A\) is invertible and
    \[
        A^{-1}=A^\conjtrans.
    \]
    Equivalently,
    \[
        A^\conjtrans A=\identmatrix{n}
        \qquad\text{and}\qquad
        AA^\conjtrans=\identmatrix{n}.
    \]
    Over \(\realnumbers\), this condition is
    \[
        A^{-1}=A^\transpose.
    \]
\end{snippetproposition}

\begin{snippetproof}{matrix-linear-isometry-criterion-proof}{matrix-linear-isometry-criterion}{Matrix criterion for linear isometries}
    Let \(\mathcal{E}\) be the canonical \orthonormalbasis of
    \(\mathbb{F}^n\). Since
    \[
        M(L_A,\mathcal{E},\mathcal{E})=A,
    \]
    the \matrix of \(L_A^\ast\) with respect to \(\mathcal{E}\) is
    \(A^\conjtrans\). By the characterization of
    \linearisometry[linear isometries],
    \[
        L_A \text{ is a \linearisometry}
        \qquad\Longleftrightarrow\qquad
        L_A^{-1}=L_A^\ast.
    \]
    In \matrix[matrices], this is exactly \(A^{-1}=A^\conjtrans\). Multiplying by
    \(A\) gives the two equivalent identities with \(\identmatrix{n}\).
\end{snippetproof}

\begin{snippetdefinition}{orthogonal-group-definition}{Orthogonal group}
    For \(n\geq 1\), the \emph{orthogonal group} is
    \[
        \operatorname{O}(n)
        =
        \left\{
            A\in\matrices_{n\times n}(\realnumbers)
            \suchthat
            A^\transpose A=\identmatrix{n}
        \right\}.
    \]
    Equivalently,
    \[
        \operatorname{O}(n)
        =
        \left\{
            A\in\matrices_{n\times n}(\realnumbers)
            \suchthat
            A^{-1}=A^\transpose
        \right\}.
    \]
    Its elements are exactly the \orthogonalmatrix[orthogonal matrices],
    equivalently the \matrix[matrices] whose associated maps are
    \linearisometry[linear isometries] of \(\realnumbers^n\) for the
    canonical \innerproduct.
\end{snippetdefinition}

\begin{snippetdefinition}{special-orthogonal-group-definition}{Special orthogonal group}
    For \(n\geq 1\), the \emph{special orthogonal group} is
    \[
        \operatorname{SO}(n)
        =
        \left\{
            A\in\operatorname{O}(n)
            \suchthat
            \det(A)=1
        \right\}.
    \]
    Its elements are the \orientationpreserving real
    \linearisometry[linear isometries] of \(\realnumbers^n\) in \matrix
    form.
\end{snippetdefinition}

\begin{snippetproposition}{orthogonal-group-product-closure}{Product of orthogonal matrices}
    Let \(n\geq 1\). If \(A,B\in\operatorname{O}(n)\), then
    \(AB\in\operatorname{O}(n)\).
\end{snippetproposition}

\begin{snippetproof}{orthogonal-group-product-closure-proof}{orthogonal-group-product-closure}{Product of orthogonal matrices}
    Since \(A^{-1}=A^\transpose\) and \(B^{-1}=B^\transpose\),
    \[
        (AB)^{-1}=B^{-1}A^{-1}=B^\transpose A^\transpose=(AB)^\transpose.
    \]
    Therefore \(AB\in\operatorname{O}(n)\).
\end{snippetproof}

\begin{snippetdefinition}{unitary-group-definition}{Unitary group}
    For \(n\geq 1\), the \emph{unitary group} is
    \[
        \operatorname{U}(n)
        =
        \left\{
            A\in\matrices_{n\times n}(\complexnumbers)
            \suchthat
            A^\conjtrans A=\identmatrix{n}
        \right\}.
    \]
    Equivalently,
    \[
        \operatorname{U}(n)
        =
        \left\{
            A\in\matrices_{n\times n}(\complexnumbers)
            \suchthat
            A^{-1}=A^\conjtrans
        \right\}.
    \]
    Its elements are exactly the \unitarymatrix[unitary matrices],
    equivalently the \matrix[matrices] whose associated maps are
    \linearisometry[linear isometries] of \(\complexnumbers^n\) for the
    canonical \innerproduct.
\end{snippetdefinition}

\begin{snippetdefinition}{special-unitary-group-definition}{Special unitary group}
    For \(n\geq 1\), the \emph{special unitary group} is
    \[
        \operatorname{SU}(n)
        =
        \left\{
            A\in\operatorname{U}(n)
            \suchthat
            \det(A)=1
        \right\}.
    \]
\end{snippetdefinition}

\end{document}
